% !TeX root = ../nc-quadric-surface-avatar.tex
% 2026-11-23 / v4.21 / cubic chart の貼合せと全域のGIT比較。

\subsection{二つのcubic chartの貼合せと重み付き逆行列写像}
\label{subsec:qa-cubic-inversion-gluing}

$\mathbb W=\mathbb P(1,2,2,3)$ の座標を
$[e_1:e_2:d:e_3]$ とし，前項の geometric 開を
$\Omega=\Omega_{\mathrm{geo}}=\mathbb W\setminus(C_0\cup C_1)$ と書く。
この式で用いる補集合は underlying closed set である。
二つの高次式が定める閉部分をその radical と取り違えない。
$\Omega^\times=\Omega\cap D_+(d)$ と置く。

\begin{proposition}[重み付きCremona対合とcubic chartの貼合せ]
\label{prop:qa-cubic-weighted-inversion-gluing}
逆行列操作 $T\mapsto T^{-1}$ が $\Omega^\times$ に誘導する対合は
\[
 \iota[e_1:e_2:d:e_3]
       =[e_3:d^2e_2:d^3:d^4e_1]
 \tag{INV}\label{eq:qa-cubic-weighted-cremona}
\]
である。右辺の四式はそれぞれ重み $3,6,6,9$ を持つ。
同じ式は $\mathbb W\setminus C_0$ から $\mathbb W$ への正則射を定め，
特に $\Omega$ 全体で定義される。しかし
\[
 \iota(\Omega\cap V(d))=\{q_0\},
 \qquad q_0=[1:0:0:0]\in C_0\setminus C_1
\]
であり，$\Omega^\times$ 上の対合は $\Omega\to\Omega$ の射へ延長できない。

全 ordered good moduli の cubic 開は，二つの $\Omega$ を
$\Omega^\times$ 上のこの対合で貼り合わせたスキームである：
\[
 M^{\mathrm{rul}}\setminus\overline N^{\mathrm{rul}}
   \simeq\Omega_f\mathop{\cup}_{\iota}\Omega_g.
 \tag{GL}\label{eq:qa-cubic-two-git-chart-gluing}
\]
ここで $\overline N^{\mathrm{rul}}\simeq\mathbb A^1_\chi$ は noncubic 閉像である。
ruling 交換は二つの chart を同じ同次座標のまま交換する。
全 $M^{\mathrm{rul}}$ 上のこの正則対合は
$\Omega^\times$ 上で式~\eqref{eq:qa-cubic-weighted-cremona} に一致し，
$\overline N^{\mathrm{rul}}$ 上では恒等である。
\end{proposition}
\begin{proof}
$K_T=JT^{\mathrm{t}}JT$ とすると，$K_{T^{-1}}$ は $K_T^{-1}$ と相似である。
$\det K_T=d^2$ だから，逆行列の基本対称式は
\[
 e_1(T^{-1})=\frac{e_3}{d^2},\qquad
 e_2(T^{-1})=\frac{e_2}{d^2},\qquad
 d(T^{-1})=\frac1d,\qquad
 e_3(T^{-1})=\frac{e_1}{d^2}.
\]
重み $(1,2,2,3)$ のスカラーとして $d^2$ を掛ければ
式~\eqref{eq:qa-cubic-weighted-cremona} を得る。
二回適用した四座標は
\[
 [d^4e_1:d^8e_2:d^9:d^{12}e_3]
\]
であり，これは重み付きスカラー $d^4$ によって元の点と同じである。
逆行列による geometric 条件の保存は
補題~\ref{lem:qa-cubic-adjugate-two-invariants} の ruling 二平面の議論による。

式の四座標が同時に零となる台は $d=e_3=0$，すなわち $C_0$ である。
従ってその補で正則射を与える。
$d=0$ では geometric 条件は $e_3\ne0$ なので，像は $q_0$ となる。
$\Omega$ は整であり，$\mathbb W$ は separated である。
仮に $\Omega\to\Omega$ への延長があれば，その $\mathbb W$ への合成は
稠密開 $\Omega^\times$ で上の正則射と一致するので，全体で一致する。
しかし $q_0\notin\Omega$ であり矛盾する。
同じ理由で $\Omega\cap V(d)$ のどの点の近傍にも $\Omega$ 値の延長はない。

一方，$f,g$ がともに可逆な square では
$f=1$ の chart の行列は $T=f^{-1}g$，
$g=1$ の chart の行列は $T'=g^{-1}f=T^{-1}$ である。
命題~\ref{prop:qa-global-cubic-four-qubit-good-moduli} の相対 chart 同型と，
命題~\ref{prop:qa-global-closed-boundary-good-images} の飽和性を用いると，
二つの good moduli chart の共通部分は正に $\Omega^\times$ となり，
その貼合せ写像が $\iota$ である。
二つの cubic 開の補は noncubic 閉像だから，表示した貼合せを得る。

ruling 交換は有限表示 square 上で $f,g$ と対応する四つの係数束を交換する。
従って二つの chart の間では同じ座標値を持つ点へ移り，
共通部分の一方の chart で見れば $\iota$ となる。
その作用は半不変式環上の作用，または good moduli の普遍性により
全 $M^{\mathrm{rul}}$ へ降下する。
補題~\ref{lem:qa-noncubic-ruling-sign-pushforward-zero} は
ordered および unordered の noncubic good moduli をともに
$\mathbb A^1_\chi$ とし，その間の射を恒等とする。
よって ruling 交換はこの閉部分の上でも scheme として恒等である。
\end{proof}

\subsection{重み二の慣性を含めた固定部分と粗二重射}

\begin{proposition}[可逆cubic領域の粗環と固定部分]
\label{prop:qa-cubic-inversion-weight-two-fixed-locus}
$D_+(d)\subset\mathbb W$ は
\[
 \operatorname{Spec}
 k[b,u,v,w]/(uw-v^2),
 \qquad
 b=\frac{e_2}{d},\quad
 u=\frac{e_1^2}{d},\quad
 v=\frac{e_1e_3}{d^2},\quad
 w=\frac{e_3^2}{d^3}
 \tag{AC}\label{eq:qa-cubic-invertible-affine-ring}
\]
であり，$\iota$ は $u,w$ を交換し $b,v$ を固定する。
その scheme-theoretic fixed locus は
\[
 u=w,\qquad (u-v)(u+v)=0
\]
で，被約な二成分である。
同次座標では $e_3=\pm de_1$ という二つの被約部分の和となる。
両成分の交差は $u=v=w=0$ であり，
$\Omega^\times$ 内では $b\ne2$ の直線である。
特に $e_1=e_3=0$ の重み二の点を除外する必要はない。

商の affine 環は
\[
 \bigl(k[b,u,v,w]/(uw-v^2)\bigr)^{\langle\iota\rangle}
       =k[b,s,v],\qquad s=u+w.
\]
$z=u-w$ と置けば粗二重射の方程式は
\[
 z^2=s^2-4v^2=(s-2v)(s+2v).
 \tag{DB}\label{eq:qa-cubic-inversion-double-cover}
\]
従って標的は smooth であり，二つの branch 因子は横断する Cartier 因子である。
源の二つの固定成分は，交差する重み二の locus ではそれぞれ Cartier ではない。
\end{proposition}
\begin{proof}
$d$ の平方根を用いて重み付き座標を $d=1$ に規格化すると，
残る $\mu_2$ は三座標 $(a,b,c)$ に
\[
 (a,b,c)\longmapsto(-a,b,-c)
\]
と作用する。ここで形式的には
$a=e_1/\sqrt d$，$b=e_2/d$，$c=e_3/d^{3/2}$ である。
この $\mu_2$ の不変式は $u=a^2,v=ac,w=c^2$ で，唯一の関係は $uw=v^2$。
これは degree-zero localization の計算でもあり，
式~\eqref{eq:qa-cubic-invertible-affine-ring} は粗スキーム上での同型である。
式~\eqref{eq:qa-cubic-weighted-cremona} はこの被覆上で $a,c$ の交換となる。
従って粗環では $u,w$ を交換し，固定イデアルは $(u-w)$ である。
商を取ると $(u-v)(u+v)=0$ となり，被約二成分の主張を得る。

$e_1=e_3=0$ では $u=v=w=0$ である。
この直線上で前項の二式は
$p_{\mathrm{adj}}=(2-b)/4$，$q_{\mathrm{adj}}=0$ となるから，
geometric 開に残る条件は $b\ne2$ である。

環を $s=u+w,z=u-w$ で書き直すと
\[
 k[b,u,v,w]/(uw-v^2)
   =k[b,s,v,z]/(z^2-s^2+4v^2).
\]
対合は $z\mapsto-z$ なので不変式環と二重射の式が従う。
$s-2v,s+2v$ の微分は交差上で一次独立だから，二つの branch 因子は横断する。
源の個々の固定成分のイデアルは
$(u-w,u-v)$ と $(u-w,u+v)$ である。
交差上ではこれらのイデアルの最小生成元は二つで，
height 一の主イデアルではない。従って個々の成分は Cartier でない。
\end{proof}

\begin{corollary}[両積可逆なunordered良モジュライの明示開]
\label{cor:qa-both-invertible-good-moduli-affine-open}
$M^{44}\subset M$ を両方の積が階数四となる飽和開の good moduli とする。
このとき
\[
 M^{44}\simeq
 \mathbb A^3_{b,s,v}\setminus
 V(s-2v,\ v-4b+8).
 \tag{M44}\label{eq:qa-both-invertible-a3-minus-line}
\]
右辺で除いているものは一本の affine 直線である。
\end{corollary}
\begin{proof}
$D_+(d)$ では $C_0$ は空であり，$C_1$ の粗方程式は
\[
 u=v=w=4b-8
\]
となる。これは $\iota$ によって点ごとに固定される曲線であり，
その商像は $s=2v,\ v=4b-8$ である。
有限群の商の fiber は幾何点の軌道なので，この曲線は位相的に飽和し，
その補の商は表示した直線の補となる。
ただし直線の scheme としての逆像は，上の二重射の式から $z^2=0$ の厚みを持つ．
ここで使った飽和性は台についてのものである．
全 $M^{\mathrm{rul}}$ から $M$ への射は ruling 交換による有限全射の商である。
実際，構成に用いた $H$-半不変式 $h_i$ の共通 chart 上で
\[
 \bigl(k[X][h_i^{-1}]^G\bigr)^{S_2}
     =k[X][h_i^{-1}]^H
\]
となり，有限群の不変式環への拡大は有限整拡大である。
この三つの affine chart は両 good moduli を覆うので，有限全射性が従う。
両積可逆な開の ordered chart は $\Omega^\times$ である。
よって上の quotient 計算が正に $M^{44}$ を与える。
\end{proof}

\subsection{随伴tensorによる全域の既存GIT商への射}

\begin{theorem}[二つの全域OU射]
\label{thm:qa-global-two-ou-morphisms}
全 ordered good moduli から二つの正則射
\[
 \rho_f,\rho_g:M^{\mathrm{rul}}\longrightarrow\mathbb W
\]
が存在し，次の性質を持つ。
\begin{enumerate}
\item $\rho_f$ は $f$ が可逆な cubic chart を $\Omega$ と同型に写し，
      $g$ が可逆な chart では式~\eqref{eq:qa-cubic-weighted-cremona} で与えられる。
      $\rho_g$ には $f,g$ を交換した記述が成立する。
\item 集合としての像はともに $\Omega\cup\{q_0\}$ であり，
      $\rho_f^{-1}(\Omega)=M_f^{\mathrm{rul}}$，
      $\rho_g^{-1}(\Omega)=M_g^{\mathrm{rul}}$ が開部分として成立する。
      $\rho_f$ は $\mathfrak D_f$ の全閉像を $q_0$ に送り，
      $\rho_g$ は $\mathfrak D_g$ の全閉像を $q_0$ に送る。
      特に noncubic 閉曲線 $\overline N^{\mathrm{rul}}\simeq\mathbb A^1_\chi$ は
      どちらの射でも $q_0$ へ写る。
\item ruling 交換 $\sigma$ に対し
      $\rho_g=\rho_f\circ\sigma$ である。
\end{enumerate}
これらの射の像は稠密な真部分であり，後の系で示す通り射は proper でない。
また二つの射の積は noncubic 閉曲線を復元できない。
\end{theorem}
\begin{proof}
全 $V_{\mathrm{cross}}$ 上で
\[
 T_f=\operatorname{adj}(f)g:
      D\otimes B\longrightarrow C\otimes A\otimes L_f
\]
を用いる。通常の四因子の二対縮約としてその Gram 不変式を取り，
\[
 E_i^{(f)}=e_i(JT_f^{\mathrm{t}}JT_f)\quad(i=1,2,3),
 \qquad D^{(f)}=\det T_f
\]
と書く。
基底変更による二つの Segre 形式の similitude の比は
\[
 \chi_f^2\frac{ac}{bd}=\zeta_f
\]
である。従って四式
$E_1^{(f)},E_2^{(f)},D^{(f)},E_3^{(f)}$ の指標は
\[
 \zeta_f,\quad\zeta_f^2,\quad\zeta_f^2,\quad\zeta_f^3
\]
で，通常の $f,g$ 座標の総次数は $8,16,16,24$ である。

この四式が同時に零にならないことを確認する。
$f$ が可逆なら
$T_f=(\det f)f^{-1}g$ は geometric tensor の非零スカラー倍であり，
前節の GIT 比較により semistable である。
従ってその四つの基本不変式は同時に零にならない。

$\operatorname{rk}f=3$ の場合には
$\operatorname{adj}(f)=\kappa\lambda^{\mathrm{t}}$ と書け，
$T_f=\kappa\psi^{\mathrm{t}}$，$\psi^{\mathrm{t}}=\lambda^{\mathrm{t}}g$ である。
$V_{\mathrm{cross}}$ の条件と階数一テンソルの交差縮約の補題から，
$\kappa,\psi$ はともに Segre 二次形式に関して非等方である。
従って
\[
 E_1^{(f)}=(\kappa^{\mathrm{t}}J\kappa)(\psi^{\mathrm{t}}J\psi)\ne0,
 \qquad E_2^{(f)}=D^{(f)}=E_3^{(f)}=0.
\]
全 $V_{\mathrm{cross}}$ では $\operatorname{rk}f\ge3$ なので，これで全場合を尽くす。

以上の四式は，実際の正則射
\[
 V_{\mathrm{cross}}\longrightarrow\mathbb W,
 \qquad
 (f,g)\longmapsto
 [E_1^{(f)}:E_2^{(f)}:D^{(f)}:E_3^{(f)}]
\]
を定める。指標が重み $(1,2,2,3)$ に従っているため，この射は $G$-不変であり，
$\mathfrak H_2^{\mathrm{rul}}$ から $\mathbb W$ への射となる。
ここで指標は $\zeta_f$ であって $\chi_0$ ではない。
$R_G(\chi_0)$ への同次数環準同型を仮定するのではなく，
good moduli space の代数空間を標的とする普遍性
\cite[Theorem 6.6]{AlperGoodModuli2013} によって
この射を $\rho_f:M^{\mathrm{rul}}\to\mathbb W$ へ降下させる。
$g$ 側にも同じ構成を行う。

$f=1$ chart では定義がそのまま $T=f^{-1}g$ の四不変式であり，
$\rho_f$ は既存の $\Omega$ との chart 同型である。
$g=1$ chart では $T_f=\operatorname{adj}(g^{-1}f)$ であるため，
$f$ も可逆な稠密開では逆行列写像に一致する。
両辺は $\Omega$ から separated な $\mathbb W$ への正則射なので，
全 $g=1$ chart で同じ式となる。
階数三の場合の計算は，$\mathfrak D_f$ 全体の像が $q_0$ であることも示す。
この閉部分は被約であり，その good moduli 閉像も被約であるから，
閉像上でも射は定値射である。

飽和性により，$M^{\mathrm{rul}}$ の各点は
$f$ 可逆な chart または $\mathfrak D_f$ の閉像に属する。
前者の像は $\Omega$，後者の像は $q_0\notin\Omega$ なので，
集合としての像と開逆像の主張を得る。
noncubic 閉部分は両階数三であり，二つの射でともに $q_0$ へ写る。
最後に，ruling 交換は $T_f,T_g$ の構成を交換するから
$\rho_g=\rho_f\circ\sigma$ である。
\end{proof}

\begin{corollary}[全good moduliの非proper性]
\label{cor:qa-global-good-moduli-nonproper}
$\rho_f,\rho_g$ は proper でなく，$M^{\mathrm{rul}}$ および $M$ は
$k$ 上 proper でない。
また $(\rho_f,\rho_g):M^{\mathrm{rul}}\to\mathbb W\times\mathbb W$ は
noncubic 閉曲線全体を $(q_0,q_0)$ へ送るため，閉埋込みではない。
\end{corollary}
\begin{proof}
$\Omega\cup\{q_0\}$ は $\mathbb W$ の稠密な真部分であり，閉じていない。
Proper 射の像は閉じているから，$\rho_f,\rho_g$ は proper でない。
$M^{\mathrm{rul}}$ と $M$ には noncubic good moduli
$\mathbb A^1_\chi$ が閉埋め込まれている。
Proper なスキームの閉部分は proper であるから，両者は $k$ 上 proper でない。
最後の主張は，閉埋込みは幾何点で単射であることと，
$\overline N^{\mathrm{rul}}\simeq\mathbb A^1_\chi$ 上で二つの射がともに定値であることから従う。
\end{proof}

\begin{remark}[既存分岐式との一致]
$\Omega^\times$ における二つの固定成分の式は
$e_3\pm de_1=0$ である。
Newton の恒等式によって，既存の OU 座標では
$D_\pm=6(e_3\pm de_1)$ となる
（系~\ref{cor:qa-ou-coarse-branch-invariant-equations}）。
上の計算はその式を $e_1=e_3=0$ の重み二の locus まで含めて
粗スキーム上で検証し，その点での源の二次錐と標的の横断 branch を区別する。
一方，$\Omega$ の補曲線 $C_0,C_1$ は非 geometric tensor の閉集合であり，
全 $H_2$ の noncubic 閉曲線ではない。
全域射 $\rho_f,\rho_g$ はこの区別を保ちながら，後者を一点 $q_0$ に送る。

再現コード \texttt{verify\_cubic\_inversion\_global\_maps.py} は，
重み付き逆写像の二乗，非対角行列の Gram 逆行列式，
階数一 tensor の普遍不変式，粗環の二重被覆式と指標を厳密算術で検算する。
射の存在，貼合せおよび降下の証明は，上記の環と相対商の議論による。
\end{remark}
